\documentclass{ximera}
\title{Inclined Cycloids}

\newcommand{\pskip}{\vskip 0.1 in}

\begin{document}
\begin{abstract}
Inclinded cycloids.
\end{abstract}
\maketitle


\section{Linear Combinations of Vectors, Part 1}
First a different way of looking at something we already know.

The basis of our entire course is the definition of the cosine and sine functions:
\[
     \cos\theta = x/r  \hskip 0.3in \sin\theta=y/r ,
\]
where $\theta$ is the polar angle of the point $(x,y)\neq 0$ and $r=\sqrt{x^2+y^2}$. Of course, the polar angle is measured counterclockwise from the positive $x$-axis.

From these defintions it follows immediately that
\[
   (x,y) = ( r\cos \theta , r\sin\theta ).
\]

The change in perspective is to think of vectors instead of coordinates and write
\[
    \langle x , y \rangle = \cos\theta \left \langle r,0 \right\rangle + \sin\theta \left\langle 0 ,r\right\rangle ,
\]
where $\langle x,y \rangle$ is the position of the point $(x,y)$ relative to the origin.

\begin{question} \label{Q3rmn8lzz}
\begin{enumerate}
\item For the non-zero vector $\langle x, y \rangle$ and $r=\sqrt{x^2+y^2}$, let $\phi$ be the angle from the vector $\langle 0,r\rangle$ to the vector $\langle x, y \rangle$, measured counterclockwise. The same geometry shows that
\[
   \langle x , y \rangle = \cos\phi \left \langle 0,r \right\rangle  + \sin\phi \left\langle  \answer{-r},0  \right\rangle.
\] 
\item Let $\alpha$ be the angle from the vector $\langle 0,r \rangle$ to the vector $\langle x, y \rangle$, measured clockwise. Then
\[
   \langle x , y \rangle = \cos\alpha \left \langle 0,r \right\rangle  + \sin\alpha \left\langle \answer{r},0  \right\rangle.
\] 
\item Let $\beta$ be the angle from the vector $\langle 0,-r\rangle$ to the vector $\langle x, y \rangle$, measured clockwise. Then
\[
   \langle x , y \rangle = \cos\beta \left \langle 0,-r \right\rangle  + \sin\beta \left\langle  \answer{-r},0  \right\rangle.
\] 
\end{enumerate}
\end{question}





\section{The Cycloid}

\begin{question}  \label{Exhjjjjnbbbbbr}

A wheel of radius $r=3$ meters rolls without slipping on the $x$-axis as a pebble $P$ attached to the wheel traces out a curve (drag the slider $u$ in the demonstration below). The wheel's axle $C$ moves  at a constant speed of $v=6$ meters/sec. Our aim is to parameterize the motion of the pebble.This means to express the $x$ and $y$ coordinates of the pebble as functions of time (measured in seconds). We assume the pebble starts at the origin at time $t=0$. 


%\href{https://www.geogebra.org/classic/vxbjccpa}{Cycloid}

 
\begin{onlineOnly}
    \begin{center}
\geogebra{vxbjccpa}{900}{600}
\end{center}
\end{onlineOnly}


\pdfOnly{Access the desmos workhseet at \href{https://www.desmos.com/calculator/57fg1na7yl}{Cycloid 25}}.

\begin{onlineOnly}
    \begin{center}
\desmos{57fg1na7yl}{900}{600}
\end{center}
\end{onlineOnly}

\href{https://www.desmos.com/calculator/57fg1na7yl}{Cycloid 25}

Here are the steps to express the coordinates (in meters) of the pebble $P$ in terms of the number of seconds since it left the origin.

\begin{enumerate}
\item First express the components of the vector $\overrightarrow{OQ}$ from the origin to the point where the wheel makes contact with the $x$-axis in terms of the rotation angle $\phi$ of the wheel. Take $\phi=0$ when the pebble is at the origin $O$.

The key idea for this is that since the wheel rolls without slipping, the length of segment $\overrightarrow{OQ}$ is equal to the length of the arc $QP$ (both shaded orange) on the wheel. Since the wheel has radius $r=3$ meters,
\[
  \overrightarrow{OQ} = \langle  \answer{3\phi} , \answer{0} \rangle. 
\]

\item Next, express the position of the circle's center $C$ relative to the origin in terms of the rotation angle $\phi$.
\[
        \overrightarrow{OC} =  \langle  \answer{3\phi} , \answer{r} \rangle. 
\]


\item Now use the ideas of Section 1 to express the components of the vector $\overrightarrow{CP}$ in terms of the vectors ${\bf v}=\langle 0,-r \rangle$, ${\bf u}=\langle -r,0 \rangle $, and the rotation angle $\phi$.

\item Activate the folder \emph{position vectors} in Line 14. Then express the position of $P$ relative to the origin (write this as a vector) in terms of the vectors $\overrightarrow{OC}$ and $\overrightarrow{CP}$.

\item Now express the components of the vector $\overrightarrow{OP}$ in terms of $\phi$.

\item Find a function 
\[
   \phi = h(t)  \, , \, t\geq 0, 
\]
that expresses the rotation angle $\phi$ in terms of the number of seconds since the pebble left the origin. Keep in mind the center of the wheel moves at a constant speed of $v=6$ m/sec.

\item Express the components of $\overrightarrow{OP}$ in terms of $t$.

\item Find coordinate functions
\[
    x=f(t)   \hskip 0.5in \text{and} \hskip 0.5in y=g(t) \, , \, t\geq 0 ,
\]
that express the coordinates of $P$ in terms of the number of seconds since the pebble left the origin.

\item Repeat part (h) in general, where the wheel has radius $r$ meters and its center moves with constant speed $v$ m/sec.
\end{enumerate}
\end{question}


\section{Rotating Vectors in the Plane}

We will have need in the next section and in later chapters to rotate a vector in the plane through $\pi/2$ radians either clockwise or counterclockwise. One way is to use geometry. Another is to think about slopes.

\begin{question} \label{Qodf3r3r}
Let 
\[
     {\bf u} = \langle 5, 2 \rangle.
\]
\begin{enumerate}
\item Find the components of the vector ${\bf v}$ obtained by rotating ${\bf v}$ counterclockwise through an angle of $\pi/2$ radians. Draw a picture

\item Find the components of the vector ${\bf w}$ obtained by rotating ${\bf v}$ clockwise through an angle of $\pi/2$ radians. Draw a picture.

\item Repeat parts (a) and (b) for the vector ${\bf u} = \langle a,b \rangle$.
\end{enumerate}
\end{question}

\begin{question} \label{Q4nvnef3ef}
\begin{enumerate} 
\item Let $\phi$ be the angle from the vector  
\[
     {\bf u} = \langle 5, 2 \rangle.
\]
to the vector $\langle {\bf p}= x, y\rangle$, measured counterclockwise. Suppose $|{\bf p}| = |{\bf u}|$ and express the components of ${\bf p}$ in terms of $\phi$.

\item Repeat part (a) but assume the angle $\phi$ is measured clockwise instead.
\end{enumerate}
\end{question}

\section{Inclined Cycloids}

\href{https://www.geogebra.org/classic/gchvgn68}{Cycloid}

 
\begin{onlineOnly}
    \begin{center}
\geogebra{gchvgn68}{900}{600}
\end{center}
\end{onlineOnly}

\begin{question} \label{Qevcefr3r430}
This is a variation of Question 2. The given information is all the same, except that the wheel rolls on a ramp inclined at the angle $\alpha$ to the ground as shown in the worksheet below.

The questions are the same, with the goal being to parameterize the pebble's motion.


\begin{onlineOnly}
    \begin{center}
\desmos{bbfxfk47vu}{900}{600}
\end{center}
\end{onlineOnly}

\href{https://www.desmos.com/calculator/bbfxfk47vu}{Inclinded Cycloid 25}


\pdfOnly{Access the desmos worksheet at \href{https://www.desmos.com/calculator/bbfxfk47vu}{142: Inclined Cycloid 25}}


\end{question}

\section{The Velocity Vector}
\emph{Velocity} is the rate of change of position with respect to time. Since position is relative, so is velocity. 

There are two types of velocity. Average velocity and instantaneous vecocity. You calculate average velocity as you would any average rate of change, here as the change in position divided by the change in time.

Calculating instantaneous velocity usually requires ideas of differential calculus. But in the special case of uniform circular motion, we do not need calculus.

The key ideas are these. At any moment, the instantaneous velocity of a point mass points in the direction of motion and has length equal to the speed of the mass.  

\begin{question} \label{Qdf3er3rbnnn}
A beetle crawls counterclockwise around a circle centered at the origin at a constant speed of $12$ cm/sec. Find the velocity of the beetle as it passes the point $(-8,10)$.
\end{question}




\end{document}